Este Trabajo de Fin de Grado consiste en un acercamiento al problema de separación de fuentes en una obra musical que persigue, dada una pieza musical, obtener mediante diferentes métodos los elementos sonoros de los que se compone.

Con este problema en mente se ha elegido el aprendizaje automático como método para llevar a cabo la tarea y se comparan diferentes métodos para valorar numéricamente el aprendizaje realizado por el sistema en una misma arquitectura. Se obtiene finalmente así una comparativa de aprendizaje, un sistema que permite obtener elementos sonoros de una pieza musical, y una interfaz gráfica que proporciona accesibilidad a un usuario promedio ajeno al tema. 